\section{A real-zero bound at prime levels and sharp interpolation at 71}

Let $p\equiv23\pmod{24}$ be prime and $m=(p+1)/24$.
Let $F$ be a nonzero cusp form with real coefficients, of weight one on
$\Gamma_0(p)$ with quadratic character $\bigl(\frac{-p}{\cdot}\bigr)$.
Assume $F(i/(py))=\sqrt p\,y F(iy)$ for all $y>0$, the positive-sign
Fricke law. This covers the theta eigenforms from nonreal ideal class
characters of $\mathbb Q(\sqrt{-p})$. The corresponding completed
Mellin transform is $\Lambda_F(s)=(\sqrt p/(2\pi))^s\Gamma(s)L(s,F)$,
interpreted by continuation.

\paragraph{Theorem.} For each integer $k\ge0$, the function
$x\mapsto\Lambda_F^{(2k)}(1/2+\sqrt x)$ has at most $m-1$ zeros on
$[0,\infty)$ counted with multiplicity. This entire function of $x$ is
defined by evenness. Equivalently, the $2k$-th derivative has at most
$2(m-1)$ zeros counted with multiplicity on the full real $s$-axis.
At $p=71$, prescribe any two points of $[0,\infty)$ in the $x$-coordinate,
allowing repetition. A unique real linear combination up to nonzero
scalar of the three completed nonabelian class-character factors realizes
exactly that multiset of zeros for $x\mapsto\Lambda_F^{(2k)}(1/2+\sqrt{x})$.
A zero of order $r$ at $x=0$ corresponds to a central zero of order $2r$.

Wei Wang's \emph{Modular $L$-functions and Chebyshev systems} gives the
corresponding geometric and Mellin-transform method and Hermite evaluation
bounds in modular weight at least two
(\href{https://arxiv.org/abs/2611.02933}{Sections 2 and 3, especially Theorem~3.3(ii)}).
Here the eta divisor supplies the geometric bound in weight one. Multiplication
of the integration kernel by $u^{2k}$ gives the same count for every fixed
even derivative.

\paragraph{Divisor bound.}
The eta-quotient criterion applied to $E(z)=\eta(z)\eta(pz)$
uses $p+1\equiv0\pmod{24}$ for both congruences. It gives the same
weight and character as $F$. The eta transformation gives the same Fricke law.
The quotient $F/E$ descends to the compact Fricke quotient $X_0(p)^+$.
The product $E$ never vanishes in the upper half-plane and starts with $q^m$.
There are two cusps on $X_0(p)$, exchanged by Fricke, so $F/E$ has only
one possible pole on the quotient, of order at most $m-1$. It has at most
that many zeros if nonconstant. A nonzero constant quotient has none.

The arc $iy$, $y>1/\sqrt p$, injects into this quotient. Indeed, for
$\left(\begin{smallmatrix}a&b\\c&d\end{smallmatrix}\right)\in\Gamma_0(p)$,
a purely imaginary image of $iy$ requires $ac y^2+bd=0$.
If both terms are nonzero then $ad$ and $bc$ are nonzero and of opposite
signs, incompatible with $ad-bc=1$. Otherwise both products vanish, so the
matrix is diagonal or antidiagonal. The first case is the identity in
$\mathrm{PSL}_2$ and the second is impossible since $p\mid c$.
In the other Fricke coset one first replaces $y$ by $1/(py)$, below the
end of the open arc. Thus $F(iy)$ has at most $m-1$ sign changes on this
open arc.

\paragraph{Passing to the Mellin integral.}
For $u=\log(\sqrt p\,y)>0$ put $K(u)=e^{u/2}F(ie^u/\sqrt p)$.
The centered derivative is $2\int_0^\infty u^{2k}K(u)\cosh(tu)\,du$.
The cusp expansion gives double-exponential decay in $u$.
For $0\le t_1<\cdots<t_l$ and $0<u_1<\cdots<u_l$ one has
$\det[\cosh(t_i u_j)]_{i,j=1}^l>0$.
Indeed, expand each entry using $(t_i^2)^h(u_j^2)^h/(2h)!$,
initially truncating at $h=M$. Cauchy--Binet expresses the determinant as a
sum of products of two generalized Vandermonde determinants with positive
weights. Each product is nonnegative by the alternant formula for
Schur polynomials, and the term using $h=0,\ldots,l-1$ is strictly positive.
This lower bound persists as $M\to\infty$.

The determinant remains strictly positive for repeated nodes in $x_i=t_i^2$
when the rows at a node of multiplicity $r$ are replaced in order by
$\partial_x^j\cosh(\sqrt{x}u)/j!$, $0\le j<r$. To see this, approximate
each repeated node by an increasing cluster of distinct, nonnegative $x$'s,
and divide the determinant by the product of differences within each
cluster. Taylor expansion in $x$ gives the displayed derivative rows in
the limit, including at $x=0$. The lower bound supplied by
$h=0,\ldots,l-1$ contains the full Vandermonde in the $x$'s. After division,
all surviving factors between different clusters have strictly positive
limits. Hence the confluent determinant is also strictly positive.

Let $u_1<\cdots<u_b$ be the sign-change points of $K$. If the transform in
$x$ had $b+1$ zeros counted with multiplicity, take the corresponding rows
of the confluent determinant, with columns at $u_1,\ldots,u_b,u$ and last
point variable. The strict sign of an ordered minor shows that this
determinant alternates sign exactly when $u$ crosses the $b$ points.
Its product with $u^{2k}K(u)$ has constant sign and nonzero integral.
Expansion in the last column gives a linear combination of vanishing
Mellin integrals and their $x$-derivatives, a contradiction.
Differentiation under the integral holds locally uniformly in $x$ by the
double-exponential decay. This also excludes an identically zero transform.

\paragraph{Sharpness at level 71.}
The class group has order seven; its six nontrivial characters yield
three inverse pairs. The three theta eigenforms are linearly independent,
as their $T_2$ eigenvalues are the three distinct numbers
$2\cos(2j\pi/7)$, $1\le j\le3$. Their positive Fricke laws follow from
Poisson summation for their binary theta components.
Berkovich and Patane give the three forms explicitly in
\href{https://arxiv.org/pdf/1302.2359}{Section~4, equations (4.2)--(4.4) and Table~6}.
Two homogeneous linear conditions on the three real coefficients admit
a nonzero solution. For a repeated point the conditions prescribe
vanishing of the value and its first $x$-derivative. The bound $m-1=2$
excludes every additional zero or multiplicity. If the solution space
had dimension at least two, imposing vanishing at another point would
still leave a nonzero solution. This proves uniqueness.

\paragraph{An explicit fourth-order central zero.}
Order the level-$71$ factors as in Berkovich and Patane's equations
(4.2)--(4.4); their $T_2$ eigenvalues in order are
$2\cos(6\pi/7),2\cos(2\pi/7),2\cos(4\pi/7)$.
Let $A=(\Lambda_j(1/2))_{j=1}^3$ and
$B=(\Lambda_j''(1/2))_{j=1}^3$. An exact projective coefficient vector
is $w=A\times B$, and the corresponding completion is
\[
 \mathcal L(s)=\det\begin{pmatrix}
    A_1&A_2&A_3\\ B_1&B_2&B_3\\
    \Lambda_1(s)&\Lambda_2(s)&\Lambda_3(s)
 \end{pmatrix}.
\]
For $k=0$ the interpolation theorem at the repeated node $x=0$ says
that the conditions $w\cdot A=w\cdot B=0$ have a one-dimensional solution
space. Thus $A\times B\ne0$. This combination has vanishing zeroth
and second central derivatives; all odd central derivatives vanish by
symmetry. The zero count proves $\mathcal L^{(4)}(1/2)\ne0$ and excludes every
other real zero of $\mathcal L$. The central cancellation occurs between
the three distinct Artin factors. For orientation, numerical quadrature
with \texttt{example71.py} gives the projective coordinates approximately
$(-0.3172807331,-0.6479315312,1)$. The determinant and zero count provide the
proof independently of these floating-point coordinates.
